% =====================================================================
% Section 1: Introduction
% Outline (Wes):
%   Intro
%     - Continual learning considers a general stream of experience.
%     - Question: how do we measure performance on arbitrary experience?
%     - Background on standard regret notions (stochastic, adversarial, ...).
%     - These don't work in general -- we need something new.
%     - We propose "perceived regret" as a general notion; give the definition.
%   Content roadmap (rest of paper)
%     - Formal definitions and assumptions.
%     - Assumptions A1, A2: upper / lower bound on the true value.
%     - Assumption A3: the examiner is learning. Why we need this
%       (vacuous upper bounds).
%     - Local notion of perceived regret w.r.t. a smaller policy class.
%     - No free lunch.
%     - Limited capacity learner and examiner.
%     - Exploration and intrinsic motivation.
%     - Practical usage and experiments.
% =====================================================================

\section{Introduction}
\label{sec:intro}

Reinforcement learning grades an agent from outside, against an optimum $V^{*}$ that a model, a prior, or a fixed comparator supplies \citep{lai1985asymptotically,jaksch2010ucrl2}. This is sound when the world is small enough to model. Continual reinforcement learning (CRL) assumes the opposite \citep{ring1994continual,abel2023definition}: under the big-world hypothesis the agent is far smaller than its environment, cannot represent or retain a fixed optimal solution, and never stops learning \citep{javed2025bigworld,lewandowski2025bigger,dohare2024plasticity}. Formally, two failures follow. At run time the agent has no estimator of $V^{*}$ that is measurable with respect to its own experience, and even given $V^{*}$ the quantity $V^{*}-f^{\pi_t}_t$ is the wrong objective, since a capacity-bounded agent should be compared to the best policy it can itself represent rather than to one it can never express \citep{simon1955behavioral,russell1995bounded,lieder2020resource,abel2024dogmas}.

Our position is that a continual agent's only honest measure of progress is one it computes from its own stream, and that to be trustworthy that measure must be produced by a module that out-learns the policy it grades. We make the measure explicit. Alongside the acting learner $\mathcal{L}$, an internal examiner $\mathcal{E}$ outputs an aspiration $\Vhi$ for the value it believes attainable and a guarantee $\Vlow$ for the value it is sure of, and the \emph{examiner regret} (also \emph{perceived regret}) is the gap $\rho_t=\Vhi-\Vlow\ge 0$. We make four formal claims. The self-report certifies nothing on its own (Proposition~\ref{prop:nothing}). A soundness condition makes $\rho_t$ an upper bound on the regret the agent can reduce (Lemma~\ref{lem:bridge}). In a coverable world this bound recovers the classical confidence-bound theory exactly (Proposition~\ref{prop:bandit}). And in a big world the bound cannot be driven to zero by any procedure, because the missing comparator is non-identifiable from the agent's stream (Lemma~\ref{lem:indist}, Theorem~\ref{thm:nofree}). The consequence is a design principle: the examiner must out-learn the learner, since estimating $\Vhi$ requires predicting the value of policies never executed. We claim no new rate; a bandit runs as the worked example throughout.

\paragraph{Where the external comparator fails.}
The case for an internal measure is that the external comparator is, in the CRL regime, either undefined or non-estimable. Three cases make this precise. First, when the world's behaviour set grows without bound, the global optimum $f^{*}_t=\sup_{\pi}f^{\pi}_t$ is a moving and possibly unattained target, so static regret against a single $V^{*}$ has no fixed referent; our within-capacity comparator $f^{K_L}_t$ stays well defined precisely because the policy class $\PiK$ is capacity-bounded with $|\PiK|\le 2^{K_L}$, which is one reason to prefer it. Second, even for a fixed action set, $f^{K_L}_t$ can depend on the unknown law $\mu$ through regions the agent need never observe, so it need not be measurable with respect to the agent's history and may admit no estimator from the stream; we make this quantitative in Lemma~\ref{lem:indist}. Third, under non-stationarity the comparator $f^{K_L}_t$ varies in $t$, so a fixed-comparator regret can be negative or vacuous, and dynamic regret \citep{wei2021master,morrill2021hindsight} reinstates a moving optimum that is again non-estimable in a big world. In all three the obstruction is that the comparator lies outside the agent's filtration. The examiner regret of Definition~\ref{def:exreg} is, by construction, a function of the agent's own stream; the rest of the paper asks what it is worth.
